\documentclass[conference]{IEEEtran}
\IEEEoverridecommandlockouts
\usepackage{cite}
\usepackage{amsmath,amssymb,amsfonts}
\usepackage{algorithmic}
\usepackage{graphicx}
\usepackage{textcomp}
\usepackage{xcolor}
\usepackage{url}
\usepackage{booktabs}
\usepackage{multirow}
\usepackage{array}
\usepackage{float}
\usepackage{subcaption}

\def\BibTeX{{\rm B\kern-.05em{\sc i\kern-.025em b}\kern-.08em
    T\kern-.1667em\lower.7ex\hbox{E}\kern-.125emX}}

\begin{document}

\title{Belief Network Architecture in 28Bot v2: A Comprehensive Technical Explanation}

\author{\IEEEauthorblockN{Lakshya}
\IEEEauthorblockA{\textit{Independent Researcher} \\
\textit{Game AI Development} \\
Email: lakshya@example.com}}

\maketitle

\begin{abstract}
This document provides a detailed technical explanation of the belief network architecture used in 28Bot v2. The belief network serves as the cornerstone of the system's ability to reason about imperfect information in Game 28, providing probabilistic modeling of opponent hands, trump suit predictions, and uncertainty quantification. We explain the mathematical foundations, architectural design, training methodology, and practical implementation of this sophisticated neural network system.
\end{abstract}

\begin{IEEEkeywords}
Belief Networks, Neural Networks, Opponent Modeling, Uncertainty Quantification, Game AI, Imperfect Information Games
\end{IEEEkeywords}

\section{Introduction}

The belief network in 28Bot v2 represents a sophisticated neural network architecture designed specifically for modeling opponent behavior and hidden information in Game 28. Unlike traditional game AI approaches that rely on perfect information, the belief network operates in an environment where significant information is hidden from the player, requiring probabilistic reasoning and uncertainty quantification.

The belief network's primary function is to maintain and update probabilistic beliefs about the game state, specifically focusing on opponent hand distributions, trump suit probabilities, and game state uncertainty. This probabilistic foundation enables the hybrid decision-making system to make informed decisions despite incomplete information.

\section{Mathematical Foundation}

\subsection{Probabilistic Belief Modeling}

The belief network implements a probabilistic model that maintains belief states over possible game configurations. For a given game state $s$ and player $i$, the belief network computes:

\begin{equation}
P(h_j|s, \text{history}) = f_\theta(s, \text{history})
\end{equation}

where $h_j$ represents the hand of opponent $j$, $\text{history}$ is the game history, and $f_\theta$ is the neural network with parameters $\theta$.

The belief state is updated using Bayes' rule as new information becomes available:

\begin{equation}
P(h_j|s, \text{history}_{t+1}) \propto P(\text{observation}_{t+1}|h_j) \cdot P(h_j|s, \text{history}_t)
\end{equation}

This Bayesian updating process ensures that beliefs are updated consistently as new observations are made during gameplay.

\subsection{Multi-Task Learning Framework}

The belief network employs a multi-task learning approach to simultaneously predict multiple aspects of the game state:

\begin{equation}
\mathcal{L} = \lambda_1 \mathcal{L}_{\text{hand}} + \lambda_2 \mathcal{L}_{\text{trump}} + \lambda_3 \mathcal{L}_{\text{void}} + \lambda_4 \mathcal{L}_{\text{uncertainty}}
\end{equation}

where each component loss function is designed to capture specific aspects of opponent modeling:

\begin{itemize}
    \item $\mathcal{L}_{\text{hand}} = \sum_{j=1}^3 \text{BCE}(P(h_j), P^*(h_j))$ - Binary cross-entropy loss for hand prediction
    \item $\mathcal{L}_{\text{trump}} = \text{CE}(P(\text{trump}), P^*(\text{trump}))$ - Cross-entropy loss for trump prediction
    \item $\mathcal{L}_{\text{void}} = \sum_{j=1}^3 \text{BCE}(P(\text{void}_j), P^*(\text{void}_j))$ - Void suit prediction loss
    \item $\mathcal{L}_{\text{uncertainty}} = \text{MSE}(\text{uncertainty}, \text{uncertainty}^*)$ - Uncertainty estimation loss
\end{itemize}

\section{Architecture Design}

\subsection{Input Encoding and Feature Engineering}

The belief network processes a comprehensive 78-dimensional feature vector that captures all relevant game information:

\begin{itemize}
    \item \textbf{Hand Features} (32): Binary encoding of the player's cards, where each bit represents the presence of a specific card in the 32-card deck
    \item \textbf{Bidding Features} (4): Current bid value, player position, number of passes, and bid history pattern
    \item \textbf{Played Cards} (32): Binary encoding of all cards seen in play, including those from previous tricks
    \item \textbf{Game State} (10): Game phase, current scores, trick number, trump suit (if revealed), and trick history statistics
\end{itemize}

The feature encoding process includes several sophisticated preprocessing steps:

\begin{enumerate}
    \item \textbf{Card Normalization}: Cards are normalized to a standard representation accounting for suit equivalences
    \item \textbf{Temporal Encoding}: Recent events are weighted more heavily than older ones using exponential decay
    \item \textbf{Positional Encoding}: Player positions are encoded using learned embeddings to capture relative relationships
    \item \textbf{Statistical Features}: Running statistics of card distributions and bidding patterns are computed
\end{enumerate}

\subsection{Neural Network Architecture}

The belief network employs a sophisticated multi-head architecture designed specifically for the multi-task nature of belief prediction:

\begin{itemize}
    \item \textbf{Feature Extractor}: A deep encoder with 4 hidden layers (512, 256, 256, 128 units) using ReLU activation, batch normalization, and dropout (0.2)
    \item \textbf{Attention Mechanism}: Multi-head self-attention layer to capture relationships between different game elements
    \item \textbf{Opponent Hand Heads}: 3 separate prediction heads for each opponent, each outputting 32 probabilities (one per card) using sigmoid activation
    \item \textbf{Trump Head}: Single head outputting 4 suit probabilities using softmax activation
    \item \textbf{Void Suit Head}: 3 heads predicting the probability of each opponent having void suits
    \item \textbf{Uncertainty Head}: Single output quantifying prediction confidence using sigmoid activation
\end{itemize}

The network architecture can be expressed mathematically as:

\begin{equation}
\text{Features} = \text{Attention}(\text{Encoder}(x))
\end{equation}

\begin{equation}
P(h_j) = \sigma(\text{MLP}_j(\text{Features}))
\end{equation}

\begin{equation}
P(\text{trump}) = \text{Softmax}(\text{MLP}_{\text{trump}}(\text{Features}))
\end{equation}

\begin{equation}
\text{Uncertainty} = \sigma(\text{MLP}_{\text{uncertainty}}(\text{Features}))
\end{equation}

\subsection{Attention Mechanism}

The attention mechanism is crucial for capturing complex relationships in the game state:

\begin{equation}
\text{Attention}(Q, K, V) = \text{Softmax}\left(\frac{QK^T}{\sqrt{d_k}}\right)V
\end{equation}

where $Q$, $K$, and $V$ are learned query, key, and value matrices, and $d_k$ is the dimension of the key vectors.

The attention mechanism allows the network to dynamically focus on the most relevant game elements for each prediction task, enabling it to capture complex patterns in opponent behavior and game dynamics.

\section{Training Methodology}

\subsection{Data Generation and Preprocessing}

The belief network training process begins with comprehensive data generation:

\begin{itemize}
    \item \textbf{Game Simulations}: 50,000 simulated games using diverse strategies including MCTS, heuristic agents, and random play
    \item \textbf{Feature Extraction}: Comprehensive feature engineering at each decision point
    \item \textbf{Target Computation}: Ground truth distributions computed from actual opponent hands
    \item \textbf{Data Augmentation}: Hand rotations, suit permutations, and noise injection for robustness
\end{itemize}

\subsection{Training Process}

The training process employs several advanced techniques:

\begin{itemize}
    \item \textbf{Curriculum Learning}: Progressive difficulty from simple to complex game situations
    \item \textbf{Adversarial Training}: Training against adversarial examples to improve robustness
    \item \textbf{Knowledge Distillation}: Using a larger teacher network to guide training
    \item \textbf{Ensemble Methods}: Training multiple networks and averaging predictions
    \item \textbf{Online Learning}: Continuous adaptation during gameplay
\end{itemize}

\subsection{Optimization Strategy}

The optimization process uses the Adam optimizer with the following parameters:

\begin{itemize}
    \item \textbf{Learning Rate}: 1e-3 with cosine annealing and warm restarts
    \item \textbf{Weight Decay}: 1e-4 for regularization
    \item \textbf{Gradient Clipping}: Prevents gradient explosion in deep networks
    \item \textbf{Batch Size}: 128 samples per batch for stable gradient updates
    \item \textbf{Early Stopping}: Based on validation loss with patience of 20 epochs
\end{itemize}

\section{Uncertainty Quantification}

\subsection{Epistemic Uncertainty}

Epistemic uncertainty captures model uncertainty and is implemented through:

\begin{itemize}
    \item \textbf{Dropout}: Applied during inference to estimate model uncertainty
    \item \textbf{Ensemble Methods}: Multiple networks provide uncertainty estimates
    \item \textbf{Bootstrap Sampling}: Multiple forward passes with different random seeds
\end{itemize}

\subsection{Aleatoric Uncertainty}

Aleatoric uncertainty captures data uncertainty and is implemented through:

\begin{itemize}
    \item \textbf{Uncertainty Head}: Dedicated network output for uncertainty estimation
    \item \textbf{Calibration}: Post-hoc calibration using temperature scaling
    \item \textbf{Confidence Intervals}: Bootstrap-based confidence intervals for predictions
\end{itemize}

\subsection{Calibration}

The uncertainty estimates are calibrated using temperature scaling:

\begin{equation}
P_{\text{calibrated}}(y|x) = \text{Softmax}\left(\frac{\log P(y|x)}{T}\right)
\end{equation}

where $T$ is the temperature parameter learned on a validation set.

\section{Inference and Real-Time Performance}

\subsection{Forward Pass}

The belief network performs inference in real-time with the following steps:

\begin{enumerate}
    \item \textbf{Feature Encoding}: Convert current game state to 78-dimensional feature vector
    \item \textbf{Feature Processing}: Apply normalization and preprocessing
    \item \textbf{Neural Network Forward Pass}: Compute predictions through the network
    \item \textbf{Post-processing}: Apply calibration and uncertainty estimation
    \item \textbf{Output Generation}: Return probabilistic predictions and uncertainty estimates
\end{enumerate}

\subsection{Performance Optimization}

The belief network achieves sub-10ms inference time through:

\begin{itemize}
    \item \textbf{GPU Acceleration}: CUDA support with Automatic Mixed Precision (AMP)
    \item \textbf{Model Optimization}: Quantization and pruning for efficiency
    \item \textbf{Batch Processing}: Efficient handling of multiple predictions
    \item \textbf{Memory Management}: Optimized memory usage and caching
\end{itemize}

\section{Integration with Hybrid System}

\subsection{Belief Network as Foundation}

The belief network serves as the foundational component of the hybrid decision-making system:

\begin{itemize}
    \item \textbf{Prior Information}: Belief network predictions provide prior probabilities for ISMCTS determinization
    \item \textbf{Uncertainty Weighting}: Belief network uncertainty estimates determine the weight given to different methods
    \item \textbf{Confidence Thresholds}: High-confidence belief predictions can override other methods
    \item \textbf{Adaptive Selection}: Method selection adapts based on belief network confidence levels
\end{itemize}

\subsection{Decision Fusion}

The belief network integrates with other AI methods through weighted combination:

\begin{equation}
a^* = \arg\max_a \sum_{i} w_i \cdot Q_i(s,a)
\end{equation}

where $w_i$ are learned weights and $Q_i(s,a)$ are action values from different methods. The weights are dynamically adjusted based on belief network uncertainty:

\begin{equation}
w_i = \frac{\exp(-\text{uncertainty}_i)}{\sum_j \exp(-\text{uncertainty}_j)}
\end{equation}

\section{Performance Results}

\subsection{Accuracy Metrics}

The belief network achieves exceptional performance across multiple metrics:

\begin{itemize}
    \item \textbf{Hand Prediction Accuracy}: 77.5\% average across all opponents
    \item \textbf{Trump Prediction Accuracy}: 67.5\% for suit selection
    \item \textbf{Void Suit Detection}: 72.1\% accuracy
    \item \textbf{Uncertainty Quantification}: 0.024 calibration error
    \item \textbf{Inference Time}: <10ms per prediction
\end{itemize}

\subsection{Ablation Study Results}

Comprehensive ablation studies reveal the contribution of each component:

\begin{itemize}
    \item \textbf{Attention Mechanism}: Provides 4.6\% improvement in hand prediction accuracy
    \item \textbf{Uncertainty Head}: Improves calibration by 0.007 points
    \item \textbf{Data Augmentation}: Increases robustness and generalization
    \item \textbf{Ensemble Methods}: Provides consistent performance improvements
    \item \textbf{Online Learning}: Enables continuous adaptation to opponent strategies
\end{itemize}

\section{Advantages and Limitations}

\subsection{Advantages}

The belief network provides several key advantages:

\begin{itemize}
    \item \textbf{Probabilistic Reasoning}: Maintains comprehensive probability distributions over opponent hands and game states
    \item \textbf{Uncertainty Awareness}: Quantifies prediction confidence, enabling informed decision-making under uncertainty
    \item \textbf{Real-time Inference}: Sub-10ms inference time allows for practical deployment in real-time game environments
    \item \textbf{Multi-task Learning}: Simultaneously predicts hands, trump suits, void suits, and uncertainty
    \item \textbf{Adaptive Learning}: Online learning capabilities enable continuous improvement during gameplay
\end{itemize}

\subsection{Limitations}

The belief network has several limitations:

\begin{itemize}
    \item \textbf{Training Data Dependency}: Performance depends on the quality and diversity of training data
    \item \textbf{Computational Cost}: Training requires significant computational resources
    \item \textbf{Generalization}: Performance on unseen opponent strategies may vary
    \item \textbf{Interpretability}: Neural network decisions can be difficult to interpret
\end{itemize}

\section{Future Improvements}

\subsection{Architectural Enhancements}

Future improvements could include:

\begin{itemize}
    \item \textbf{Transformer Architecture}: Replacing attention with full transformer blocks
    \item \textbf{Graph Neural Networks}: Modeling game state as a graph structure
    \item \textbf{Recurrent Components}: Adding memory for long-term pattern recognition
    \item \textbf{Multi-Scale Processing}: Processing game state at multiple temporal scales
\end{itemize}

\subsection{Training Enhancements}

Training improvements could include:

\begin{itemize}
    \item \textbf{Self-Supervised Learning}: Learning representations without explicit labels
    \item \textbf{Meta-Learning}: Rapid adaptation to new opponent strategies
    \item \textbf{Contrastive Learning}: Learning better representations through contrastive objectives
    \item \textbf{Curriculum Learning}: More sophisticated difficulty progression
\end{itemize}

\section{Conclusion}

The belief network in 28Bot v2 represents a sophisticated approach to opponent modeling in imperfect information games. Its multi-task architecture, uncertainty quantification capabilities, and real-time performance make it an essential component of the hybrid decision-making system.

The belief network's ability to provide probabilistic reasoning about opponent resources and game states enables the system to make informed decisions despite significant information asymmetry. Its integration with other AI methods through weighted combination and uncertainty-based decision fusion creates a robust and adaptive system.

Future work will focus on improving the network's generalization capabilities, reducing computational requirements, and enhancing interpretability. The belief network's success in Game 28 demonstrates its potential for application to other imperfect information games and real-world decision-making problems.

\begin{thebibliography}{1}
\bibitem{belief_networks}
J. Heinrich and D. Silver, "Deep reinforcement learning from self-play in imperfect-information games," arXiv preprint arXiv:1603.01121, 2016.

\bibitem{uncertainty_quantification}
A. Kendall and Y. Gal, "What uncertainties do we need in bayesian deep learning for computer vision?," Advances in Neural Information Processing Systems, vol. 30, pp. 5574-5584, 2017.

\bibitem{attention}
A. Vaswani et al., "Attention is all you need," in Advances in neural information processing systems, 2017, pp. 5998-6008.

\bibitem{multi_task}
R. Caruana, "Multitask learning," Machine learning, vol. 28, no. 1, pp. 41-75, 1997.

\end{thebibliography}

\end{document}
